\documentclass[11pt]{article}
\usepackage[margin=1in]{geometry}
\usepackage[T1]{fontenc}
\usepackage[utf8]{inputenc}
\usepackage{amsmath,amssymb,amsthm}
\usepackage{enumitem}

\title{ICPC World Finals 2022\\V. Three Kinds of Dice}
\author{}
\date{}

\begin{document}
\maketitle

\section*{Problem Summary}

We are given two dice $D_1$ and $D_2$, where $D_1$ has an advantage over $D_2$. We want two extremal
values over all possible third dice $D_3$:
\begin{itemize}[leftmargin=*]
    \item the minimum possible $\text{score}(D_3,D_2)$ subject to $\text{score}(D_3,D_1)\ge 1/2$;
    \item the maximum possible $\text{score}(D_3,D_1)$ subject to $\text{score}(D_3,D_2)\le 1/2$.
\end{itemize}

\section*{Key Observations}

\begin{itemize}[leftmargin=*]
    \item Fix one face value $v$ of $D_3$. Its contribution against $D_1$ and $D_2$ is the point
    \[
    p(v)=\bigl(\text{score}(D_1,v),\ \text{score}(D_2,v)\bigr),
    \]
    where $\text{score}(D_i,v)$ is the expected score of $D_i$ against a deterministic one-face die
    showing $v$.
    \item A general die $D_3$ is just a multiset of face values, so the pair
    \[
    \bigl(\text{score}(D_1,D_3),\ \text{score}(D_2,D_3)\bigr)
    \]
    is a convex combination of points $p(v)$. Thus the set of all achievable pairs is exactly the convex
    hull of the candidate points.
    \item The point $p(v)$ changes only when $v$ crosses a face value of $D_1$ or $D_2$. So it is enough
    to consider $1$, $\max+1$, and for every face value $x$ appearing on either die also $x$ and $x-1$.
    \item Once the convex hull is known, the required extrema are intersections of that hull with the
    half-planes $x \le 1/2$ and $y \ge 1/2$.
\end{itemize}

\section*{Algorithm}

\begin{enumerate}[leftmargin=*]
    \item Sort both dice. If necessary, swap them so that the first one is the stronger die.
    \item Enumerate all candidate face values $v$ described above.
    \item For each $v$, compute $p(v)$ by binary searching in the two sorted dice.
    \item Sort these points by $x$ and keep, for each equal $x$, only the largest $y$.
    \item Build the upper convex hull of the remaining points.
    \item Scan the hull:
    \begin{itemize}[leftmargin=*]
        \item maximize $y$ over the part with $x \le 1/2$;
        \item minimize $x$ over the part with $y \ge 1/2$;
        \item if an edge crosses one of those threshold lines, interpolate on that edge.
    \end{itemize}
    \item Convert back using $\text{score}(D_3,D_i)=1-\text{score}(D_i,D_3)$.
\end{enumerate}

\section*{Correctness Proof}

We prove that the algorithm returns the correct answer.

\paragraph{Lemma 1.}
For every possible die $D_3$, the pair
\[
\bigl(\text{score}(D_1,D_3),\ \text{score}(D_2,D_3)\bigr)
\]
lies in the convex hull of the candidate points.

\paragraph{Proof.}
If the faces of $D_3$ are $v_1,\dots,v_m$, then
\[
\text{score}(D_i,D_3)=\frac{1}{m}\sum_{j=1}^{m}\text{score}(D_i,v_j)
\]
for $i=1,2$. Hence the pair for $D_3$ is the average of the points $p(v_j)$, so it lies in the convex
hull. \qed

\paragraph{Lemma 2.}
Every point in the convex hull of the candidate points is achievable by some die $D_3$.

\paragraph{Proof.}
The candidate points themselves are produced by deterministic one-face dice. Any convex combination
with rational coefficients can be realized by taking the corresponding number of copies of each face
value. Therefore the whole convex hull is achievable. \qed

\paragraph{Lemma 3.}
It is sufficient to consider only the candidate face values listed by the algorithm.

\paragraph{Proof.}
For a fixed die $D_i$, the value $\text{score}(D_i,v)$ depends only on how many faces of $D_i$ are less
than $v$ and how many are equal to $v$. These counts change only when $v$ crosses a face value already
present on $D_1$ or $D_2$. Thus every open interval between consecutive critical values yields the same
point, so one representative per interval is enough. \qed

\paragraph{Theorem.}
The algorithm outputs the correct two extremal scores.

\paragraph{Proof.}
By Lemmas 1 and 2, the achievable region for
\[
\bigl(\text{score}(D_1,D_3),\ \text{score}(D_2,D_3)\bigr)
\]
is exactly the convex hull of the candidate points. Therefore minimizing $\text{score}(D_3,D_2)$ under
$\text{score}(D_3,D_1)\ge 1/2$ is equivalent to maximizing $\text{score}(D_2,D_3)$ under
$\text{score}(D_1,D_3)\le 1/2$, which is precisely the highest point of the hull inside $x \le 1/2$. The
second answer is symmetric.

The optimum of a linear objective over a convex polygon is attained either at a vertex or on an edge
crossing the boundary line. Thus scanning the upper hull and interpolating on crossing edges yields the
correct values. \qed

\section*{Complexity Analysis}

Let $n_1$ and $n_2$ be the numbers of faces on the two given dice. We create $O(n_1+n_2)$ candidate
values. Each candidate point is evaluated with binary searches in the sorted dice, so the total running
time is $O((n_1+n_2)\log(n_1+n_2))$. The memory usage is $O(n_1+n_2)$.

\section*{Implementation Notes}

\begin{itemize}[leftmargin=*]
    \item The implementation first ensures the first die is the stronger one, so the output formula is
    always the same.
    \item \texttt{long double} is sufficient for the required precision.
\end{itemize}

\end{document}
